\chapter{Experimental Platform and Methodology}
\label{ch:methodology}
% Prose below is the plain-English synthesis (05). Fuller material and Chinese notes:
% ../dissertation_prep/01_draft_ch3_methodology.md.

\section{Hardware and Software Environment}
\label{sec:env}
All experiments ran on ARCHER2. Each node has two AMD EPYC 7742 processors, 128 physical
cores, and eight NUMA regions of 16 cores each. The software stack was PrgEnv-gnu 8.4.0,
gcc 11.2.0, and cray-mpich 8.1.27. PETSc 3.24.4 was built from source as
\texttt{arch-omp-opt} with \texttt{--with-openmp=1 --with-openmp-kernels=1} and
\texttt{-O3}. One thread per rank is used as the MPI-only baseline.
% TODO: expand from 01 §3.1 (Core Complex Dies, shared L3, 2.25 GHz base frequency).

\section{Benchmark Problems and Drivers}
\label{sec:benchmarks}
The 2D benchmark is PETSc \texttt{ex2}, a five-point stencil discretisation of the Laplace
operator on an $m \times n$ grid. The 3D benchmark is a driver written for this project, a
seven-point stencil on an $m \times n \times p$ grid; it was added because the
surface-to-volume ratio differs from 2D. Profiling uses \texttt{ex2\_repeat}
(\Cref{ch:profiling}).

\section{Solver Configuration}
\label{sec:solver}
The default \texttt{ex2} solver is GMRES with block Jacobi. At $4500 \times 4500$ and
above, the error norm reached order $10^{10}$, so the iteration did not converge within
the default limit. The solver was changed rather than the tolerance relaxed. Production
runs use
\begin{verbatim}
-ksp_type cg -pc_type gamg -ksp_rtol 1e-5 -ksp_atol 1e-10
\end{verbatim}
CG applies because the discretised Laplacian is symmetric positive definite. GAMG is used
because multigrid convergence is close to mesh-independent for elliptic problems. Across
5.0M to 41.0M unknowns the solver takes 11--12 iterations with an error norm of order
$10^{-2}$ ($0.026$ at $4500 \times 4500$). Constant iteration counts mean runtime
differences reflect per-iteration and setup cost, not iteration count.

GAMG cost has two phases: setup (\texttt{PCSetUp}, building the coarse hierarchy) and
solve (\texttt{KSPSolve}). In a single-solve run, setup is about 45\% of runtime; in
applications the preconditioner is reused across solves. \Cref{ch:singlenode,ch:multinode}
report total time including setup; \Cref{ch:profiling} isolates the solve phase.
% TODO(G1): replace "order 10^10" with a table (default vs CG+GAMG: error/iters/runtime).

\section{Configuration Space and Selection Rules}
\label{sec:configspace}
Three rules reduce the space.

Threads per rank are restricted to $\{1,2,4,8\}$. A single-node sweep showed 16 or more
threads per rank were slower at every core count, and a single rank (pure OpenMP) reduced
runtime by under 20\% from 1 to 128 threads. At eight threads or fewer a rank stays within
one or two NUMA regions.

Nodes are fully populated: $\text{ppn} \times \text{threads} = 128$, giving four layouts
$128\times1$, $64\times2$, $32\times4$, and $16\times8$. Same node count means same core
count, so a comparison isolates rank--thread balance from core count.

A configuration at size $N$ on $C$ cores is included only if $10^4 < N/C < 10^6$. This
sets the valid node counts per size. Four sizes are used: 5.0M, 10.0M, 20.25M, and 41.0M.
The 2D and 3D grid dimensions are matched per nominal size (for example
$4500^2 \approx 273^3 \approx 20.25\,\text{M}$).

\section{Job Configuration and Thread Placement}
\label{sec:placement}
Each rank and its threads are pinned to a contiguous set of cores with:
\begin{verbatim}
OMP_PLACES=cores   OMP_PROC_BIND=close
srun --hint=nomultithread --distribution=block:block --exact
\end{verbatim}
% TODO: from 01 §3.5 — why this pinning follows from the NUMA layout.

\section{Metrics}
\label{sec:metrics}
Runtime is the maximum \texttt{Time (sec)} over ranks from \texttt{-log\_view}, taken as
the median over repetitions. Scaling is computed per layout against that layout's smallest
node count (node-relative, not single-core; this differs from the feasibility study). For
a layout $\ell$ at $n$ nodes,
\[
  S_\ell(n) = \frac{T_\ell(n_0)}{T_\ell(n)}, \qquad
  E_\ell(n) = \frac{S_\ell(n)}{n/n_0},
\]
where $n_0$ is the smallest node count run for that problem size.

\section{Reproducibility}
\label{sec:repro}
% TODO: from 01 §3.7 — scripts, drivers, raw logs, one-command figure regeneration.
